\documentclass[12pt, a4paper]{article}
\usepackage[utf8]{inputenc}
\usepackage[italian]{babel}
\usepackage{graphicx} % Per inserire i diagrammi UML

\title{Relazione Progetto: Sistema Gestione Biblioteca}
\author{Il Tuo Nome}
\date{Maggio 2026}

\begin{document}

\maketitle

\section{Introduzione}
Il presente progetto riguarda la realizzazione di un sistema per la gestione dei prestiti e della ricerca di libri in una biblioteca universitaria.

\section{Requisiti del Sistema}
\subsection{Casi d'Uso Principali}
Il sistema si concentra su due funzionalità core:
\begin{itemize}
    \item \textbf{Ricerca Libro:} Permette all'utente di verificare la presenza di un titolo nel catalogo.
    \item \textbf{Gestione Prestito:} Gestisce la transazione di prestito verificando la disponibilità del volume.
\end{itemize}

% Qui inserirai il diagramma UML che hai esportato prima
% \begin{figure}[h]
%     \centering
%     \includegraphics[width=0.8\textwidth]{casi_uso.png}
%     \caption{Diagramma dei Casi d'Uso}
% \end{figure}

\section{Architettura e Design}
Il sistema è stato sviluppato seguendo i principi SOLID, in particolare il principio di Singola Responsabilità (SRP), separando i modelli dei dati dalla logica di business.

\end{document}